% !TEX TS-program = pdflatex
% !TEX encoding = UTF-8 Unicode

% This file is a template using the "beamer" package to create slides for a talk or presentation
% - Giving a talk on some subject.
% - The talk is between 15min and 45min long.
% - Style is ornate.

% MODIFIED by Jonathan Kew, 2008-07-06
% The header comments and encoding in this file were modified for inclusion with TeXworks.
% The content is otherwise unchanged from the original distributed with the beamer package.

\documentclass{beamer}
\setbeamercovered{invisible}



% Copyright 2004 by Till Tantau <tantau@users.sourceforge.net>.
%
% In principle, this file can be redistributed and/or modified under
% the terms of the GNU Public License, version 2.
%
% However, this file is supposed to be a template to be modified
% for your own needs. For this reason, if you use this file as a
% template and not specifically distribute it as part of a another
% package/program, I grant the extra permission to freely copy and
% modify this file as you see fit and even to delete this copyright
% notice. 


\mode<presentation>
{
  \usetheme{Warsaw}
  % or ...

  % or whatever (possibly just delete it)
}


\usepackage[english]{babel}
% or whatever

\usepackage[utf8]{inputenc}
% or whatever

\usepackage{times}
\usepackage[T1]{fontenc}
% Or whatever. Note that the encoding and the font should match. If T1
% does not look nice, try deleting the line with the fontenc.

%%% MATH RELATED

\input{macros.tex}

\title{MATH 350-2 Advanced Calculus} 
\subtitle
{} % (optional)

\author[W.R. Casper] % (optional, use only with lots of authors)
{W.R. Casper}
% - Use the \inst{?} command only if the authors have different
%   affiliation.

\institute[California State University Fullerton] % (optional, but mostly needed)
{
  Department of Mathematics\\
  California State University Fullerton}
% - Use the \inst command only if there are several affiliations.
% - Keep it simple, no one is interested in your street address.

\subject{Talks}
% This is only inserted into the PDF information catalog. Can be left
% out. 



% If you have a file called "university-logo-filename.xxx", where xxx
% is a graphic format that can be processed by latex or pdflatex,
% resp., then you can add a logo as follows:

% \pgfdeclareimage[height=0.5cm]{university-logo}{university-logo-filename}
% \logo{\pgfuseimage{university-logo}}



% Delete this, if you do not want the table of contents to pop up at
% the beginning of each subsection:
\AtBeginSubsection[]
{
  \begin{frame}<beamer>{Outline}
    \tableofcontents[currentsection,currentsubsection]
  \end{frame}
}


% If you wish to uncover everything in a step-wise fashion, uncomment
% the following command: 

%\beamerdefaultoverlayspecification{<+->}


\begin{document}

\begin{frame}
  \titlepage
\end{frame}

\begin{frame}{Outline}
  \tableofcontents
  % You might wish to add the option [pausesections]
\end{frame}

% Since this a solution template for a generic talk, very little can
% be said about how it should be structured. However, the talk length
% of between 15min and 45min and the theme suggest that you stick to
% the following rules:  

% - Exactly two or three sections (other than the summary).
% - At *most* three subsections per section.
% - Talk about 30s to 2min per frame. So there should be between about
%   15 and 30 frames, all told.

\section{Real Analysis Lecture 18}

\subsection{Uniform continuity}

\begin{frame}{Uniform continuity}
\begin{defn}
Let $(S,d_S)$ and $(T,d_T)$ be metric spaces and $A\subseteq S$.
We say that a function $f: A\rightarrow T$ is \textbf{uniformly continuous} on $A$ if for all $\epsilon > 0$ there exists $\delta > 0$ such that
$$d_S(x,y)< \delta \Rightarrow d_T(f(x),f(y))< \epsilon$$
for all $x,y\in A$.
\end{defn}
\begin{itemize}
\item $f$ continuous on $A$ means $f$ is continuous at every $a\in A$
\item $\delta$ can depend on $\epsilon$ and $a$ (not on $x$)
\item with uniform continuity, $\delta$ can't depend on $a$ ...
\end{itemize}
\end{frame}

\begin{frame}{Examples}
\begin{itemize}
\item $x$ is uniformly continuous on $\mathbb R$
\item $x^2$ is NOT uniformly continuous on $\mathbb R$
\item $1/x$ is continuous on $(0,1)$
\item $1/x$ is NOT uniformly continuous on $(0,1)$
\end{itemize}
\end{frame}

\begin{frame}{Relationship with compactness}

\begin{thm}[Heine]
Let $(S,d_S)$ and $(T,d_T)$ be metric spaces and $A\subseteq S$ be compact.
Suppose that $f: A\rightarrow T$ is continuous on $A$, then $f$ is uniformly continuous on $A$.
\end{thm}
\end{frame}


\subsection{Infinite series}

\begin{frame}{Basic definition}
\begin{defn}
An \textbf{infinite series} is an ordererd pair of sequences $(\{a_n\},\{s_n\})$ where 
$$s_n = a_1+a_2+\dots+a_n.$$
The sequence $\{s_n\}$ is called the \textbf{sequence of partial sums} and $s_n$ is called the $n$'th partial sum.
We say the series converges or diverges if $\{s_n\}$ converges or diverges.
\end{defn}
\textbf{Notation:}
$$a_1+a_2+a_3+\dots\quad\text{or}\quad \sum_{n=1}^\infty a_n.$$
We write
$$\sum_{n=1}^\infty a_n = L$$
to mean that the limit of the sequence $s_n$ is $L$.
We call $L$ the \textbf{sum} of the series.
\end{frame}

\begin{frame}{Examples}
$$\sum_{n=1}^\infty \frac{1}{n^2+n} = 1.$$
\pause
Telescoping:
\pause
\begin{align*}
s_n &= \frac{1}{1^2+1} + \frac{1}{2^2+2} +  \frac{1}{3^2+3} + \dots + \frac{1}{n^2+n}\\
    &= \left(\frac{1}{1}-\frac{1}{2}\right) + \left(\frac{1}{2}-\frac{1}{3}\right) +\dots+ \left(\frac{1}{n}-\frac{1}{n+1}\right)\\
    &= 1-\frac{1}{n+1}.
\end{align*}
\end{frame}

\begin{frame}{Examples}
$$\sum_{n=1}^\infty r^n = \frac{r}{1-r},\quad |r| < 1.$$
\pause
Geometric:
\pause
$$s_n = r + r^2 + r^3 + \dots + r^n= r(1+r+r^2+\dots+r^{n-1}) = r(1+s_{n-1})$$
\pause
$$s_n = r + r^2 + r^3 + \dots + r^n= (r+r^2+\dots+r^{n-1}) + r^n = s_{n-1} + r^n$$
\pause
$$s_n = \frac{r-r^{n+1}}{1-r}$$
\end{frame}

\begin{frame}{Examples}
$$\sum_{n=1}^\infty (-1)^{n+1}\frac{1}{n} = \ln(2)$$
\pause
Taylor series:
\pause
$$\ln(1-x) = -x - \frac{1}{2}x^2 - \frac{1}{3}x^3 - \frac{1}{4}x^4 - \dots$$
\pause
Converges for $-1\leq x < 1$.
\end{frame}

\begin{frame}{Tests for Convergence}
\begin{thm}[Cauchy condition]
The series $\sum a_n$ converges if and only if for all $\epsilon > 0$ there exists $N\in\mathbb Z_+$ such that $n\geq N$ implies
$$\left\lvert a_{n} + a_{n+1} + \dots + a_{n+p}\right\rvert < \epsilon$$
for all $p\in\mathbb Z_+$.
\end{thm}
\end{frame}

\begin{frame}{Tests for Convergence}
\begin{defn}
An \textbf{alternating series} is a series of the form $\sum (-1)^{n+1}a_n$ with $a_n\geq 0$ for all $n$.
\end{defn}
\begin{thm}[Alternating Series Test]
If $a_n$ is a decreasing sequence converging to $0$, then $\sum (-1)^{n+1}a_n$ converges to a real number $s$ and
$$|s-s_n| < a_{n+1}\quad\forall\ n\in\mathbb Z_+.$$
\end{thm}
\end{frame}

\begin{frame}{Tests for Convergence}
\begin{defn}
We say $\sum a_n$ is \textbf{absolutely convergent} if $\sum |a_n|$ converges.\\
We say $\sum a_n$ is \textbf{conditionally convergent} if it is convergent but not absolutely convergent.
\end{defn}
\begin{thm}[Absolute Convergence Test]
If a series is absolutely convergent, then it is convergent.
\end{thm}
\end{frame}

\begin{frame}{Tests for Convergence}
\begin{thm}[Comparison Test]
If $a_n>0$ and $b_n > 0$ and there exists $c>0$ and $N\in\mathbb Z_+$ with $a_n < cb_n$ for all $n\geq N$ then
$$\sum b_n\ \ \text{converges}\quad \Rightarrow\quad\sum a_n\ \ \text{converges}$$
\end{thm}
\end{frame}

\begin{frame}{Tests for Convergence}
\begin{thm}[Limit Comparison Test]
If $a_n>0$ and $b_n > 0$ and suppose that $\lim a_n/b_n = c > 0$.  Then
$$\sum b_n\ \ \text{converges}\quad \Leftrightarrow\quad\sum a_n\ \ \text{converges}$$
\end{thm}
\end{frame}

\begin{frame}{Limit infimum/limit supremum}
\begin{defn}
Let $a_n$ be a sequence.  Then
$$\liminf a_n = \lim_n \inf\{a_k: k\geq n\}$$
$$\limsup a_n = \lim_n \sup\{a_k: k\geq n\}$$
\end{defn}
\begin{thm}
If $a_n$ is bounded, then $\limsup a_n$ and $\liminf a_n$ exist.
\end{thm}
More generally, we write $\liminf a_n = \pm\infty$ and $\limsup a_n = \pm \infty$.
\end{frame}

\begin{frame}{Tests for Convergence}
\begin{thm}[Ratio Test]
Let $r = \liminf |a_{n+1}/a_n|$ and $R = \limsup |a_{n+1}/a_n|$.
\begin{enumerate}[(a)]
\item if $R < 1$ then the series is absolutely convergent
\item if $r > 1$ then the series is divergent
\item if $r \leq 1\leq R$ the test is inconclusive
\end{enumerate}
\end{thm}
\end{frame}

\begin{frame}{Tests for Convergence}
\begin{thm}[Root Test]
Let $\rho = \limsup \sqrt[n]{|a_n|}$.
\begin{enumerate}[(a)]
\item if $\rho < 1$ then the series is absolutely convergent
\item if $\rho > 1$ then the series is divergent
\item if $\rho = 1$ the test is inconclusive
\end{enumerate}
\end{thm}
\end{frame}

\subsection{Infinite products}

\begin{frame}{Basic definition}
\begin{defn}
An \textbf{infinite product} is an ordererd pair of sequences $(\{a_n\},\{p_n\})$ where 
$$p_n = a_1a_2\dots a_n.$$
The sequence $\{p_n\}$ is called the \textbf{sequence of partial products} and $p_n$ is called the $n$'th partial product.
We say the series converges or diverges if $\{p_n\}$ converges or diverges.
\end{defn}
\textbf{Notation:}
$$a_1a_2a_3\dots\quad\text{or}\quad \prod_{n=1}^\infty a_n.$$
We write
$$\prod_{n=1}^\infty a_n = L$$
to mean that the limit of the sequence $p_n$ is $L$.
We call $L$ the \textbf{product} of the series.
\end{frame}

\begin{frame}{Examples}
Weierstrass product of sine
$$\sin(\pi x) = \pi z\prod_{n=1}^\infty \left(1-\frac{x^2}{n^2}\right)$$
\begin{itemize}
\item first thought of by Euler.
\item can be used to prove 
$$\sum \frac{1}{n^2} = \frac{\pi^2}{6}.$$
\end{itemize}
\end{frame}

\begin{frame}{Examples}
Let $p_1=2,p_2=3,p_3=5,\dots$ be the prime integers.\\
The infinite product
$$\zeta(s) = \prod_{n=1}^\infty \frac{1}{1-p_n^{-s}}$$
converges for $s>1$.\\
This product is called the \textbf{Riemann zeta function}.
\end{frame}

\begin{frame}{Examples}
\begin{thm}
For any real number $s>1$,
$$\zeta(s) = \sum_{n=1}^\infty \frac{1}{n^s}.$$
\end{thm}
Therefore
$$\frac{1}{(1-1/2^2)(1-1/3^2)(1-1/5^2)(1-1/7^2)(1-1/11^2)\dots} = \frac{\pi^2}{6}.$$
\end{frame}

\begin{frame}{Relation to series}
\begin{thm}
Assume $a_n > 0$ for all $n$.
Then $\prod (1+a_n)$ converges if and only if $\sum a_n$ converges.
\end{thm}
\end{frame}

\begin{frame}{Relation to series}
\begin{defn}
We say that a product $\prod (1+a_n)$ \textbf{converges absolutely} if $\prod (1+|a_n|)$ converges.
\end{defn}
\begin{thm}
If $\prod (1+a_n)$ converges absolutely (equiv. $\sum a_n$ converges absolutely), then $\prod (1+a_n)$ converges.
\end{thm}
\end{frame}

\end{document}


